\section{Conclusion and Future Directions}
\label{sec:conclusion}

This Systematic Literature Review synthesized 18 empirical and foundational studies to evaluate the effect of prototyping on identifying missing and unclear requirements within software engineering projects. Grounded in the IREB/CPRE requirements engineering paradigm and a structured tripartite analytical framework (\textit{Rationale, Operational Protocol, and Taxonomical Classification}), our findings yield four principal conclusions:

\begin{enumerate}
    \item \textbf{Modality and Purpose (RQ6.1):} Prototyping transcends simple UI illustration; under the Prototyping Assessment Model (PAM), prototypes function as cognitive boundary objects. In early elicitation, exploratory probes surface tacit domain rules; in validation, interactive simulations catalyze real-time requirement evolution and semantic disambiguation.
    \item \textbf{Comparative Efficacy (RQ6.2):} Rigorous empirical evidence from multi-cohort experimental families confirms that Paper Prototyping significantly outperforms traditional unstructured interviews and JAD in gross functional requirement discovery volume ($p = 0.02$). However, passive visual prototypes do not guarantee specification completeness ($p = 0.17$) or correctness ($p = 0.59$) for non-visual backend logic, refuting the assumption that prototyping is an automated panacea.
    \item \textbf{Requirement Categorization (RQ6.3):} Visual prototypes excel predominantly in eliciting functional workflows and UI structures. The traditional vulnerability regarding Non-Functional Requirements (NFRs) is successfully mitigated through augmented modalities, specifically Loud Paper Prototyping (LPP) combining Wizard-of-Oz manipulation with Think-Aloud protocols ($p = 0.001$), and video-based narrative scenarios.
    \item \textbf{Risks, Metrics, and Agile Integration:} High-fidelity prototypes introduce severe risks of premature idea cementing and release pressure. In contrast, integrating low-fidelity sketches into Agile Enriched User Stories during Backlog Grooming minimizes sprint rework by 30--40\%. Furthermore, formalizing linguistic ambiguity metrics (ARR and $\Delta_{\text{interp}}$) provides a rigorous mathematical foundation to quantify requirement clarity.
\end{enumerate}

\subsection{Future Work}
As part of future research trajectories, we plan to execute the proposed quasi-experimental research blueprint across industrial development teams. This initiative will combine automated language model ambiguity classifiers with live interactive elicitation sessions, establishing an automated, evidence-based pipeline for requirements quality assurance.
